\section*{TÓM TẮT ĐỒ ÁN}
\phantomsection\addcontentsline{toc}{section}{\numberline {}TÓM TẮT ĐỒ ÁN}

Trong những năm gần đây, phẫu thuật thẩm mỹ đã trở thành một nhu cầu phổ biến nhằm cải thiện diện mạo khuôn mặt. Tuy nhiên, sự thay đổi đáng kể về mặt hình học và cấu trúc của các vùng cơ bản trên khuôn mặt sau phẫu thuật (như nâng mũi, gọt hàm, cắt mí) đặt ra một thách thức vô cùng lớn đối với các hệ thống xác thực khuôn mặt truyền thống. Hầu hết các mô hình học sâu hiện nay (như FaceNet, VGG-Face, ArcFace) học đặc trưng toàn thể khuôn mặt (global representation), do đó sự thay đổi cục bộ ở một vài vùng phẫu thuật sẽ làm sai lệch hoàn toàn vector đặc trưng tổng thể, dẫn đến sai sót nghiêm trọng trong quá trình so khớp (tăng tỷ lệ loại bỏ sai - False Rejection Rate).

Đồ án này đề xuất giải pháp tối ưu hóa quá trình so khớp khuôn mặt sau phẫu thuật thẩm mỹ dựa trên phương pháp \textbf{Tích hợp vùng cục bộ thích ứng (Adaptive Local Patch Fusion)} tích hợp trên nền tảng thư viện DeepFace. Hệ thống sử dụng bộ phát hiện mốc khuôn mặt để tự động cắt nhỏ khuôn mặt thành 3 vùng cục bộ độc lập: Vùng Mắt \& Trán (đại diện cho cấu trúc xương bất biến), Vùng Mũi, và Vùng Miệng \& Cằm. Đặc trưng vector (embedding) của từng vùng được trích xuất riêng biệt bằng mô hình CNN mà không cần chạy lại bộ phát hiện mặt. Dựa trên khoảng cách Cosine của các vùng, giải thuật sẽ tự động phân tích sự bất thường (anomaly) để phát hiện vùng bị can thiệp phẫu thuật, từ đó điều chỉnh trọng số (weights) một cách động: giảm trọng số vùng bị chỉnh sửa và tăng trọng số của các vùng bất biến. Kết quả thử nghiệm thực tế cho thấy giải pháp đạt độ chính xác cao và tăng khả năng kháng nhiễu đối với các khuôn mặt bị thay đổi cơ học sâu sắc, mở ra hướng ứng dụng thực tiễn hiệu quả trong quản lý khách hàng tại các cơ sở thẩm mỹ.

\newpage
\section*{ABSTRACT}
\phantomsection\addcontentsline{toc}{section}{\numberline {}ABSTRACT}

\begin{sloppypar}
In recent years, plastic-surgery has emerged as a popular means to enhance facial appearance. However, significant alterations in the geometry and local structure of facial features after surgery (such as rhinoplasty, jaw shaving, or double eyelid surgery) present a major challenge to traditional face verification systems. Most modern deep learning models (e.g., FaceNet, VGG-Face, ArcFace) learn holistic representations of the face; hence, local changes in surgically altered regions severely distort the global feature vector, causing high False Rejection Rates (FRR).

This thesis proposes an optimized face verification framework for post-plastic-surgery faces using \textbf{Adaptive Local Patch Fusion} built upon the DeepFace library. The system automatically segments the face into three local, independent patches based on facial coordinates: the Eyes \& Forehead patch (representing invariant bone structure), the Nose patch, and the Mouth \& Chin patch. Local embeddings are extracted individually using deep CNNs under skip-detection mode. By analyzing local Cosine distances, the algorithm dynamically detects anomalies corresponding to surgically altered regions and shifts the fusion weights away from modified patches toward invariant ones. Experimental results demonstrate that the proposed solution substantially improves verification accuracy and robustness against mechanical facial changes, providing a practical utility for security and customer management systems in aesthetic clinics.
\end{sloppypar}

\newpage
\cleardoublepage